% =====================================================================
% 致谢（可选）
%
% 以下文字与附件 3 致谢页完全一致，逐字照录（含附件 3 自带的格式说明
% 段落）。章标题照录为“致    谢”（中间空 4 格）。
%   内容中文用宋体小四号，1.5 倍行距，首行缩进 2 字符。
% =====================================================================

% 正文标题“致    谢”（中间空 4 格）；目录条目用不带空隙的“致谢”。
\ujsunnumberedchapter{致\qquad 谢}[致谢]

致谢（可选）应以简短的文字对在课题研究与设计（论文）撰写过程中直接给予帮助的指导教师、答疑教师和其他人员或单位表示自己的谢意，以示对他人劳动的尊重。语言应实事求是，切记浮夸之词。在结论后单置一页。

“致谢”二字用黑体小二号，居中，中间空4格，段前间距为2行，段后间距为1行。

内容中文用宋体小四号，1.5倍行距，首行缩进2字符。
